

\vspace{-0.25cm}
\section{Does Agent Trust Align with Human Trust?}
\label{sec:alignment}
\vspace{-1.5mm}








In this section, we aim to explore the fundamental relationship between agent and human trust, \ie, whether or not agent trust aligns with human trust. This provides important insight regarding the feasibility of utilizing LLM agents to simulate human trust behavior as well as more complex human interactions that involve trust. First, we propose a new concept \textit{behavioral alignment} and discuss its distinction from existing alignment definitions. Then, we conduct extensive studies to investigate whether or not LLM agents exhibit alignment with humans regarding trust behavior. 

\vspace{-2mm}
\subsection{Behavioral Alignment}
\vspace{-1mm}
Existing alignment definitions predominantly emphasize \textit{values} that seek to ensure the safety and helpfulness of LLMs~\citep{ji2023ai,shen2023large,wang2023aligning}, which cannot fully characterize the landscape of multifaceted alignment between LLMs and humans.
Thus, we propose a new concept of \textit{behavioral alignment} to 
characterize the LLM-human  analogy regarding \textit{behavior}, which involves both actions and the associated reasoning processes that underlie them.
Because actions evolve over time and the reasoning that underlies them involves multiple factors, we define  \textbf{\textit{behavioral alignment}}
as the analogy between LLMs and humans concerning factors impacting behavior, namely \textbf{\textit{behavioral factors}}, and action dynamics, namely \textbf{\textit{behavioral dynamics}}.

Based on the definition of behavioral alignment, we aim to answer: \textit{does agent trust align with human trust?} As for \textit{behavioral factors}, existing human studies have shown that three basic factors impact human trust behavior including reciprocity anticipation~\citep{berg1995trust,coxHowIdentifyTrust2004b}, risk perception~\citep{BOHNET2004467} and prosocial preference~\citep{ferrerTrustGames2019a}. We examine whether agent trust aligns with human trust along these three factors. Although  \textit{behavioral dynamics} vary for different humans and agent personas, we analyze whether agent trust has the same patterns across multiple turns as human trust in the Repeated Trust Game. 

Besides analyzing the trust behavior of LLM agents and humans based on quantitative measurements (\eg, the \textit{amount sent} from trustor to trustee), we also explore the use of \textit{BDI} to interpret the reasoning process with which LLM agents justify their actions, which can further validate whether LLM agents manifest an underlying reasoning process analogous to human cognition.

























\vspace{-2mm}
\subsection{Behavioral Factor 1: Reciprocity Anticipation}

Reciprocity anticipation, the expectation of a reciprocal action from the other player, can positively influence human trust behavior~\citep{berg1995trust}. The effect of reciprocity anticipation exists in the Trust Game but not in the Dictator Game (Section \ref{game:Dictator Game and trust game} Games 1 and 2) because trustee cannot return money in the Dictator Game, which is the only difference between these games.
Thus,  to determine whether LLM agents can anticipate reciprocity, we compare their behaviors in these Games.






















\begin{wrapfigure}{t}{0.54\textwidth}
\vspace{-4mm}
  \centering
    \includegraphics[width=0.54\textwidth]{images/dic_trust_plot1.pdf}
    \vspace{-6mm}
\caption{\textbf{The Comparison of Average Amount Sent  for LLM Agents and Humans in the Trust Game and the Dictator Game}. 
}
  \label{fig:dic_vs_trust}
  \vspace{-2mm}
\end{wrapfigure}

First, we analyze trust behaviors based on the average amount of money sent by  human or LLM agents.
As shown in Figure \ref{fig:dic_vs_trust}, human studies show that humans exhibit a higher level of trust in the Trust Game than in the Dictator Game ($\$6.0$ vs. $\$3.6$, $p$-value = $0.01$ using One-Tailed Independent Samples t-test)~\citep{coxHowIdentifyTrust2004b}, indicating that reciprocity anticipation enhances human trust. Similarly, GPT-4 ($\$6.9$ vs. $\$6.3$, $p$-value = $0.05$ using One-Tailed Independent Samples t-test) 
also shows a higher level of trust in the Trust Game with statistical significance, implying that reciprocity anticipation can enhance agent trust. However, LLMs with fewer parameters (\eg, Llama2-13b) do not show this tendency in their trust behaviors for the Trust and Dictator Games.




Then, we further analyze GPT-4 agents' BDI to explore whether they can anticipate reciprocity in their reasoning (the complete BDIs are in Appendix~\ref{Dictator Game vs. Trust Game}). 
Typically, in the Trust Game, one persona's BDI emphasizes ``\textit{putting faith in people}'', which implies the anticipation of the goodness of the other player, and ``\textit{reflection of trust}''. However, in the Dictator Game, one persona's BDI focuses on concepts such as ``\textit{fairness}'' and ``\textit{human kindness}'', which are not directly tied to trust or reciprocity. Thus, we can observe that GPT-4 shows distinct BDI outputs in the Trust and Dictator Games.

Based on the above analysis of the amount sent and BDI, we find that 
\textbf{GPT-4 agents exhibit human-like reciprocity anticipation in trust behavior}. Nevertheless, \textbf{LLMs with fewer parameters  (\eg, Llama2-13b) do not show an awareness of reciprocity from the other player}.














\vspace{-0.15cm}
\subsection{Behavioral Factor 2: Risk Perception}
\begin{wrapfigure}{t}{0.55\textwidth}
\vspace{-3mm}
  \centering
    \includegraphics[width=0.55\textwidth]{images/Risky_Dictator_Game_and_Trust_Game.pdf}
    \vspace{-6mm}
  \caption{\textbf{Trust Rate (\%) Curves for LLM Agents and Humans in the MAP Trust Game and the Risky Dictator Game.}  The metric Trust Rate indicates the portion of trustors opting for trust given \(p\).}
  \label{fig:Risk Perception}
  \vspace{-4mm}
\end{wrapfigure}


Existing human studies have demonstrated the strong correlation between trust behavior and risk perception, suggesting that human trust will increase as risk decreases
~\citep{hardin2002trust_risk,trust_risk_1,coleman1994foundations_risk}. We aim to explore whether LLM agents can perceive the risk associated with their trust behaviors through the MAP Trust Game and the Risky Dictator Game (Section~\ref{game:Dictator Game and trust game} Games 3 and 4), where risk is represented by the probability \((1-p)\) (defined in Section~\ref{game:Dictator Game and trust game}).




As shown in Figure~\ref{fig:Risk Perception}, we measure human trust (or agent trust) by the portion choosing to trust the other player in the whole group, namely the Trust Rate (\%).
Based on existing human studies~\citep{BOHNET2004467}, when the probability \(p\) is higher, the risk for trust behaviors is lower, and more humans choose to trust, manifesting a higher Trust Rate, which indicates that human trust rises as risk falls. Similarly, we observe a general increase in agent trust as risk decreases for LLMs including GPT-4, GPT-3.5-turbo-0613, and text-davinci-003.
In particular, we can see that the curves of humans and GPT-4 are more aligned compared with other LLMs, implying that GPT-4 agents' trust behaviors dynamically adapt to different risks in ways most aligned with humans. LLMs with fewer parameters (\eg, Vicuna-13b) do not exhibit the similar tendency of Trust Rate as the risk decreases.



We further analyze the BDI of GPT-4 agents to explore whether they can perceive risk through reasoning (complete BDIs in Appendix~\ref{BDI:map trust game}). Typically, under high risk ($p=0.1$), one persona's BDI mentions ``\textit{the risk seems potentially too great}'', suggesting a cautious attitude.  Under low risk ($p=0.9$), one persona's BDI reveals a strategy to ``\textit{build trust while acknowledging potential risks}'', indicating the willingness to engage in trust-building activities despite residual risks. Such changes in BDI reflect how GPT-4 agents perceive risk changes in the reasoning underlying their trust behaviors.

Through the analysis of Trust Rate Curves and BDI, we can infer that \textbf{GPT-4 agents manifest human-like risk perception in trust behaviors}. Nevertheless, \textbf{LLMs with fewer parameters (\eg, Vicuna-13b) often do not perceive risk changes in their trust behaviors}.







\vspace{-0.2cm}
\subsection{Behavioral Factor 3: Prosocial Preference}
\begin{wrapfigure}{t}{0.53\textwidth}
\vspace{-4mm}
  \centering
    \includegraphics[width=0.53\textwidth]{images/lottery_curve.pdf}
    \vspace{-6mm}
  \caption{\textbf{Lottery Rates (\%) for LLM Agents and Humans in the Lottery Gamble Game and the Lottery People Game}. Lottery Rate indicates the portion of choosing to gamble or trust the other player.}
  \label{fig:Prosocial Preference}
  \vspace{-4mm}
\end{wrapfigure}
\vspace{-0.5mm}
Human studies have found that the  prosocial preference, referring to humans' inclination to trust other humans in contexts involving social interaction~\citep{ferrerTrustGames2019a,fetchenhauerBetrayalAversionPrincipled2012}, also  plays a key role in human trust behavior. 
We study whether LLM agents have prosocial preference in trust behaviors by comparing their behaviors in the Lottery  Gamble Game (LGG) and the Lottery  People Game (LPG) (Section \ref{game:Dictator Game and trust game} Game 5). The only difference between these two games is the effect of prosocial preference in LPG, because the winning probability of gambling $p$ in LGG is the same as the reciprocation probability $p$ in LPG.


\vspace{-0.5mm}
As shown in Figure~\ref{fig:multi_round_res}, existing human studies have demonstrated that more humans are inclined to place trust in other humans over relying on pure chance ($54\%$ vs. $29\%$)~\citep{fetchenhauerBetrayalAversionPrincipled2012}, implying that the prosocial preference is essential for human trust.
We can observe the same tendency in most LLM agents except Vicuna-13b.  
For GPT-4 in particular, a much higher percentage of the personas choose to trust the other player over gambling ($72\%$ vs. $21\%$),  illustrating that the prosocial preference is also an important factor for GPT-4 agents' trust behaviors. 



When interacting with humans, GPT-4's BDI typically indicates a preference to ``\textit{believe in the power of trust}'', in contrast to gambling, where the emphasis shifts to ``\textit{believing in the power of calculated risks}''. The comparative analysis of reasoning processes (complete BDIs in Appendix~\ref{Lottery Game}) demonstrates that GPT-4 agents tend to embrace risk when involved in social interactions. This tendency aligns closely with the concept of prosocial preference observed in human trust behaviors.

The analysis of the Lottery Rates and BDI suggests that \textbf{LLM agents, especially GPT-4 agents,  demonstrate human-like prosocial preference in trust behaviors, except Vicuna-13b}.













\vspace{-0.2cm}
\subsection{Behavioral Dynamics}
\label{bahavior_dynamic}
\vspace{-1mm}
Besides behavioral factors, we also aim to investigate whether LLM agents align with humans regarding trust behavioral dynamics over turns in the Repeated Trust Game (Section \ref{game:Dictator Game and trust game} Game 6).

Admittedly, existing human studies show that the dynamics of human trust over turns are complex due to human diversity. The complete results from 16 groups of human experiments are shown in Appendix~\ref{appendix:Repeated Trust Game Results:Person}~\citep{jones1998experience}. We still observe three common  patterns for human trust behavioral dynamics in the Repeated Trust Game: \textbf{\textit{First}, the amount returned is usually larger than the amount sent in each round}, which is natural because the trustee will receive $\$3N$ when the trustor sends $\$N$; \textbf{\textit{Second}, the ratio between amount sent and returned generally remains stable except for the last round}. In other words, when the amount sent increases, the amount returned is also likely to increase. And when the amount sent remains unchanged, the amount returned also tends to be unchanged. This reflects the stable relationship between trust and reciprocity in humans. Specifically, the ``Returned/3$\times$Sent Ratio'' in Figure \ref{fig:multi_round_res} is considered stable if the fluctuation between successive turns is within $10\%$; \textbf{\textit{Third}, the amount sent (or returned) does not manifest frequent fluctuations across turns}, illustrating a relatively stable underlying reasoning process in humans over successive turns. Typically, Figure \ref{fig:multi_round_res} Humans (a) and (b) show these three  patterns.









\begin{wrapfigure}{t}{0.57\textwidth}
\vspace{-4mm}
  \includegraphics[width=0.57\textwidth]{images/multi_round_img/section/3.5/merge.pdf}
    \includegraphics[width=0.57\textwidth]{images/multi_round_img/section/4/merge.pdf}
    \includegraphics[width=0.57\textwidth]{images/multi_round_img/section/person/merge.pdf}
    \vspace{-5mm}
  \caption{\textbf{Results of GPT-4, GPT-3.5 and Humans in the Repeated Trust Game.} The blue lines indicate the amount sent or returned for each round. The red lines imply the ratio of the amount returned to three times of the amount sent for each round.}
  \label{fig:multi_round_res}
  \vspace{-2mm}
\end{wrapfigure}

We conducted 16 groups of the Repeated Trust Game with GPT-4 or GPT-3.5-turbo-0613-16k (GPT-3.5), respectively. For the two players in each group, the personas  differ to reflect human diversity and the LLMs are the same. Complete results are shown in the Appendix~\ref{sec:GPT-4_multi_round_result},~\ref{sec:GPT-3.5_multi_round_result} and typical examples are shown in Figure \ref{fig:multi_round_res} GPT-3.5 (a) (b) and  GPT-4 (a) (b). Then, we examine whether the aforementioned three patterns observed in human trust behavior also manifest in trust behavioral dynamics of GPT-4  (or GPT-3.5). For GPT-4 agents, we discover that these  patterns generally exist in all $16$ groups ($87.50\%$, $87.50\%$, and $100.00\%$ of all results show these three patterns, respectively).  However, fewer GPT-3.5 agents manifest these patterns ($62.50\%$, $56.25\%$, and $43.75\%$ hold these three patterns, respectively). The experiment results show that \textbf{GPT-4 agents demonstrate highly human-like patterns in their trust behavioral dynamics}. Nevertheless, \textbf{a relatively large portion of GPT-3.5 agents fail to show human-like patterns in their dynamics}, indicating such behavioral patterns may require stronger cognitive capacity.





Through the comparative analysis of LLM agents and humans in the 
\textit{behavioral factors} and \textit{dynamics} associated with trust behavior, evidenced in both
their \textit{actions} and \textit{underlying reasoning processes}, our second core finding is as follows:

\vspace{-2.5mm}
\begin{center}
\begin{tcolorbox}[width=0.93\linewidth, boxrule=3pt, colback=gray!20, colframe=gray!20]
\textbf{Finding 2:} 
GPT-4 agents exhibit high \textit{behavioral alignment} with humans regarding trust behavior under the framework of Trust Games, although other LLM agents, which possess fewer parameters and weaker capacity, show relatively lower \textit{behavioral alignment}.



\end{tcolorbox}
\end{center}
\vspace{-2.5mm}
This finding underscores the potential of using LLM agents, especially GPT-4, to simulate human trust behavior, encompassing both  \textit{actions} and  underlying \textit{reasoning processes}. This paves the way for the simulation of more complex human interactions and institutions. This finding deepens our understanding of the fundamental analogy between LLMs and humans and opens avenues for research on LLM-human alignment beyond values.














